\begin{exsol}{22}{\MR{} Analytical Problem 5.5}{165}{Exam}{B3}
\asked Only the seller is informed: she values the good at \$0 (probability $x$) or \$10, and the
buyer's value is the seller's plus a known premium. Show that with a \$10 premium a single price
of \$10 supports honest reports and trade in every state; then, with a \$5 premium, decide whether
report-contingent prices can make the seller honest while both sides are willing to trade.

\method \MR's mediator and incentive-compatibility method (Ch.~5, pp.~141--143), here with
one-sided private information: a price schedule $p(\hat s)$ depending on the seller's report
$\hat s \in \{0, 10\}$, honesty constraints for the two seller types, and participation
constraints for both seller types and for the buyer, who learns the state from an honest report.
Trade probabilities below one are the screening device of \MR's incentive-efficiency
discussion (pp.~143--145).

\solution
Write $s \in \{0, 10\}$ for the seller's value and $s + g$ for the buyer's, with $g = 10$ or
$g = 5$. Trade is efficient in both states because $g > 0$.

\textbf{(a) Premium \$10: one price of \$10 always.} Let $p(0) = p(10) = 10$ and trade always.
\begin{itemize}
  \item \emph{Honesty.} The report changes neither the price nor whether trade occurs, so the
  seller cannot gain by misreporting; honest reporting is (weakly) optimal for both types.
  \item \emph{Seller participation.} Type 0 gains $10 - 0 = 10 \ge 0$; type 10 gains
  $10 - 10 = 0$, so she is willing to sell.
  \item \emph{Buyer participation.} When $s = 0$ the good is worth \$10 to him and he pays \$10
  (gain 0); when $s = 10$ it is worth \$20 (gain 10). He is willing to buy in each state, even
  if he could infer the state, so the scheme survives ex post.
\end{itemize}
\begin{center}
\begin{tabular}{lccc}
\toprule
State & Price & Seller's gain & Buyer's gain \\
\midrule
$s = 0$ (prob.\ $x$) & 10 & 10 & 0 \\
$s = 10$ (prob.\ $1-x$) & 10 & 0 & 10 \\
\bottomrule
\end{tabular}
\end{center}
Efficient trade occurs in every state, with honest revelation. (Any constant price in
$[10,\, 20 - 10x]$ also works if the buyer only needs a non-negative \emph{expected} gain, since
$x(10-p) + (1-x)(20-p) = 20 - 10x - p$; \$10 is the price that also works state by state.)

\textbf{(b) Premium \$5: is there an honest schedule with trade in both states?} Now the buyer
is worth \$5 when $s = 0$ and \$15 when $s = 10$. Suppose trade always occurs at
$p(0)$ and $p(10)$.
\begin{enumerate}
  \item \emph{Honesty of type 0}: $p(0) \ge p(10)$ (she sells either way, so only the price
  matters).
  \item \emph{Honesty of type 10}: $p(10) - 10 \ge p(0) - 10$, i.e.\ $p(10) \ge p(0)$.
  Together: $p(0) = p(10) \equiv P$. When the report cannot change whether trade happens, it
  cannot change the price either.
  \item \emph{Participation of type 10}: $P \ge 10$.
  \item \emph{Buyer's willingness to pay}: after an honest report of $0$ the buyer knows the
  good is worth \$5 to him, so he buys only if $P \le 5$.
\end{enumerate}
Conditions 3 and 4 are incompatible: \textbf{no such prices exist}. More generally, trade in
all states with honesty forces one pooled price that must lie in $[10,\, 0 + g]$, which is
non-empty iff $g \ge 10$: the premium must be at least the gap between the seller's two values.
That is exactly why the \$10 case works and the \$5 case fails.

\textbf{What can be achieved instead.} Honesty is restored by making the claim ``\$10'' costly
through a lower chance of trade. Let trade occur for sure at $p(0) \le 5$ after report 0, and
with probability $q$ at $p(10) \ge 10$ after report 10. Type 0 stays honest iff
$p(0) \ge q\,p(10)$; with $p(0) = 5$ and $p(10) = 10$ this allows any $q \le 1/2$. Type 10 does
not mimic type 0, since selling at \$5 loses \$5; the buyer gains $5 - 5 = 0$ and
$15 - 10 = 5$ in the two states. The best such scheme trades with probability $q = 1/2$ when
the good is worth \$10 to the seller, and the expected loss against full efficiency is
$(1-x)\cdot\tfrac12\cdot 5 = 2.5\,(1-x)$.

\emph{Caveat on timing.} If the buyer had to commit to a single price before anything is
revealed and judged it by his expected gain, one price $P \in [10,\, 15 - 10x]$ would do, which is
possible iff $x \le 1/2$. Once prices depend on the report, though, the report reveals the state,
and the buyer's willingness to trade must hold state by state, as above.

\interp Private information does not need two-sided uncertainty to destroy efficient trade. With
\emph{interdependent} values the buyer's value moves with the seller's information. The low-cost
seller always wants to be paid like a high-cost one, so honesty pools the prices. The buyer will
not pay the high-cost price for a good he then knows is worth little. When the gains from trade
are large ($g \ge 10$) a single price covers everybody; when they are small, the informed party's
informational rent cannot be paid without giving up trade, and the incentive-efficient
arrangement rations trade with the high-value seller. This is the logic of Akerlof's lemons model,
which \MR{} discuss under adverse selection in the same chapter, in bargaining form.

\verify Each constraint is checked by substitution in \texttt{check\_D.py}: seller and buyer gains
in the \$10 case; the condition $10 \le P \le g$, which holds iff $g \ge 10$; and $q = 1/2$ from
$5 = q \times 10$, the buyer's gains (0 and 5) and type 10's loss ($5 - 10 < 0$) in the \$5 case;
the interim range closes at $x = 1/2$.

\answer{With a \$10 premium, a constant price of \$10 and trade always is honest (reports change
nothing) and individually rational for both seller types and for the buyer in each state. With
a \$5 premium, no: honesty with trade in both states forces $p(0) = p(10) = P$, the \$10 seller
needs $P \ge 10$, and a buyer who learns the good is worth \$5 pays at most \$5. Efficient trade
with honesty needs a premium of at least \$10. The best honest scheme trades at \$5 after report
0 and with probability $1/2$ at \$10 after report 10.}
\end{exsol}
